\section{Resultados Esperados}

De acuerdo con las brechas identificadas en el diagnóstico inicial y los objetivos
estratégicos de Machu Picchu Exclusive Tours E.I.R.L., la implementación de Odoo 19 se
orientó a lograr una transformación estructural en la gestión de sus servicios turísticos
de lujo. Se establecieron los siguientes resultados esperados:

\begin{itemize}
    \item \textbf{Centralización y seguridad de la información confidencial:} se proyectó la digitalización total de la base de datos de pasajeros (PAX), permitiendo el almacenamiento seguro de documentos de identidad (pasaportes) en un repositorio centralizado. Esto debe eliminar definitivamente el riesgo de ciberseguridad asociado al manejo de fotos de pasaportes en dispositivos personales y chats de WhatsApp.
    \item \textbf{Control operativo y trazabilidad de itinerarios:} se espera estandarizar el seguimiento de los servicios mediante un tablero visual Kanban, permitiendo un control riguroso de itinerarios que varían entre 2 a 20 días. El sistema debe garantizar que cada etapa (desde el borrador de cotización hasta la ejecución final) cuente con la documentación necesaria, como vouchers y boletos de tren o avión.
    \item \textbf{Optimización financiera de las reservas:} se buscó formalizar el registro manual de adelantos (30\% al 50\%) y saldos. El resultado esperado es contar con una visibilidad inmediata de los estados de pago para mitigar riesgos financieros antes de proceder con la compra de boletos y servicios no reembolsables.
    \item \textbf{Mejora en la eficiencia comercial y conversión:} se proyectó una reducción del 40\% en los tiempos de respuesta al cliente internacional y una disminución equivalente en errores administrativos derivados de la gestión manual. Estas optimizaciones deben traducirse en un incremento de entre el 15\% y 20\% en la tasa de conversión de cotizaciones a ventas efectivas.
    \item \textbf{Inteligencia de negocios para la expansión internacional:} se espera que el sistema genere reportes en tiempo real sobre nacionalidades con mayor demanda, rentabilidad por tour y desempeño de canales de captación. Esta información es crítica para fortalecer la presencia de la marca en los mercados norteamericano y europeo y reducir la dependencia de intermediarios.
    \item \textbf{Fidelización y atención personalizada:} mediante el uso del módulo de actividades y el historial del CRM, se busca automatizar recordatorios para seguimientos post-viaje (como tarjetas de cumpleaños o fin de año), incrementando la tasa de clientes recurrentes en un 20\%.
\end{itemize}
